\section{Experimental Results}
We present the primary evaluation results of Nexa in terms of throughput, energy efficiency, and area overhead.

\subsection{Throughput and Energy Efficiency}
Table~\ref{tab:results} summarizes the performance comparison of Nexa against the baseline architectures. Nexa achieves a peak throughput of 840.4 GSOPS (Giga Synaptic Operations Per Second) at 1.20 GHz, which represents a 6.75x speedup over the GPU baseline and a 3.42x speedup over the ASIC-based Apex accelerator. This performance improvement is primarily due to the decentralized spike routing protocol, which avoids global synchronization bottlenecks, and the look-ahead prefetching in the SWC.

\begin{table*}[t]
  \centering
  \caption{Comparison of Nexa with State-of-the-Art Neural Simulation Baselines.}
  \label{tab:results}
  \begin{tabular}{lccccc}
    \toprule
    \textbf{Architecture} & \textbf{Node} & \textbf{Area} & \textbf{Freq.} & \textbf{Throughput} & \textbf{Efficiency} \\
     & (nm) & (mm$^2$) & (GHz) & (GSOPS) & (pJ/SOP) \\
    \midrule
    GPU Baseline (Nexa-G) & 7 & N/A & 1.40 & 124.5 & 350.2 \\
    ASIC Baseline (Apex) & 16 & 12.4 & 0.80 & 245.8 & 84.1 \\
    Meridian v2 & 14 & 18.2 & 1.00 & 310.2 & 62.4 \\
    \midrule
    \textbf{Nexa Core (This Work)} & 7 & 8.5 & 1.20 & 840.4 & 12.8 \\
    \bottomrule
  \end{tabular}
\end{table*}

Furthermore, Nexa demonstrates exceptional energy efficiency. By keeping synaptic weight accesses local to the on-chip SWC, Nexa minimizes energy-expensive off-chip HBM accesses. Specifically, Nexa consumes only 12.8 pJ per Synaptic Operation (pJ/SOP), compared to 350.2 pJ/SOP for the GPU baseline and 84.1 pJ/SOP for the Apex accelerator. This translates to an energy reduction of over 27x and 6.5x, respectively.

\subsection{Cache Hit Rate and Memory Traffic}
To analyze the impact of the look-ahead prefetching in the SWC, we measure the cache hit rate across the three workloads. Figure~\ref{fig:results_chart} illustrates the relative energy consumption per spike for each workload compared to the baseline.

\begin{figure}[h]
  \centering
  \begin{tikzpicture}
    \begin{axis}[
      ybar,
      bar width=12pt,
      width=0.95\columnwidth,
      height=5.2cm,
      ylabel={Energy (nJ / Spike)},
      xlabel={Workloads},
      symbolic x coords={S1-Sparse, S2-Medium, S3-Dense},
      xtick=data,
      nodes near coords,
      nodes near coords align={vertical},
      ymin=0, ymax=12,
      legend style={at={(0.5,-0.25)}, anchor=north, legend columns=-1},
      axis x line*=bottom,
      axis y line*=left,
      enlarge x limits=0.25,
      ymajorgrids=true,
      grid style={dashed, gray!30},
      every node near coord/.append style={font=\tiny}
    ]
      \addplot[fill=iscablue!70] coordinates {(S1-Sparse, 2.1) (S2-Medium, 4.3) (S3-Dense, 8.7)};
      \addplot[fill=iscaaccent!70] coordinates {(S1-Sparse, 0.4) (S2-Medium, 0.9) (S3-Dense, 1.8)};
      \legend{Baseline, Nexa}
    \end{axis}
  \end{tikzpicture}
  \caption{Normalized energy consumption per spike across different workloads.}
  \label{fig:results_chart}
\end{figure}

For the S1-Sparse workload, Nexa achieves a cache hit rate of 94.2\%, which remains above 89.5\% even for the highly irregular S3-Dense workload. This high hit rate validates our prefetching strategy, which successfully anticipates post-synaptic accesses based on incoming routing tables.
